\subsubsection{The Origins of Adversarial Training in GANs}
\addcontentsline{toc}{subsubsection}{\protect\numberline{\thesubsubsection}{The Origins of Adversarial Training in GANs}}
Adversarial training approaches include a wide variety of architectures where two networks or loss functions compete against each other. Adversarial approaches are inspired by Generative Adversarial Networks (GANs) (Box \ref{box:gan}) \citep{goodfellow2014generative}. In the machine learning literature on causal inference, adversarial training has been applied both to trade off outcome modeling and treatment modeling tasks during representation learning, as well as to trade off estimation and regularization of IPW weights. GANs have also been used directly as generative models for counterfactual and treatment effect distributions. 

\begin{TCBBox}[label=box:gan]{Generative Adversarial Networks (GAN)}In GANs, two networks, a discriminator network $D$ and a generator network $G$, play a zero-sum game like cops and robbers. The generator network's job is to learn a distribution from which the training data $X$ could have credibly been generated. In each training batch, the generator produces a new outcome (originally images, but could be IPW weights, counterfactuals or treatment effects) by drawing a random noise sample from a known distribution $Z$ (e.g. Gaussian) and transforming it into outcomes with the function $G(Z)=\hat{X}$. The discriminator's job is to learn a function $D(X)=P(X\ is\ real)$ that can distinguish whether the outcome is from the training data $X$, or whether it is a ``fake" $\hat{X}$ created by the generator. The generator then receives a negative version of the discriminator's loss, a penalty that is proportional to how well it was able to ``deceive" the discriminator.  The discriminator's loss can be the log loss,  Jensen-Shannon divergence \citep{goodfellow2014generative}, the Wasserstein distance \citep{arjovsky2017wasserstein, gulrajani2017improved}, or any number of divergences and IPMs. Formally, the generator attempts to minimize the following loss function,
\begin{equation*}
\argmin_{G}=\underbrace{\Ex_{X}}_{\text{ real dist.}}[\underbrace{\mathcal{L}(D(X)}_{P(X\ is\ real)}] + \underbrace{\Ex_{Z}}_{\text{fake dist.}}[1-\underbrace{\mathcal{L}(D(G(Z))}_{P(\hat{X}\text{is real})}]
\end{equation*}
where the first quantity is the discriminator's estimated probability data from $X$ is indeed real, and the second quantity is the discriminator's estimate that a generated quantity from the distribution $Z$ is real. 

Because the discriminator is trying to catch the generator, its objective is to maximize the same function,
\begin{equation*}
\argmax_{D}=\underbrace{\Ex_{X}}_{\text{ real dist.}}[\underbrace{\mathcal{L}(D(X)}_{P(X \text{is real})}] + \underbrace{\Ex_{Z}}_{\text{fake dist.}}[1-\underbrace{\mathcal{L}(D(G(Z))}_{P(\hat{X}\text{is real})}]
\end{equation*}
In practice, the discriminator and the generator are trained either alternatingly or simultaneously, with the discriminator increasing its ability to discern between real and fake outcomes over time, and the generator increasing its ability to deceive the discriminator. The idea is that the adaptive loss function created by the discriminator can coax the generator out of local minima to generate superior outcomes. Results by these models have been impressive, and many of the fake portraits and ``deepfake" videos circulating online in recent years are generated by this architectures. The advantage of GANs is that they can impressively learn very complex generative distributions with limited modeling assumptions. The disadvantage of GANs is that they are difficult and unreliable to train, often plateauing in local optima.
\end{TCBBox}
\label{Box:GANs}

\subsubsection{GANs as Generative Models of Treatment Effect Distributions (GANITE)}
\textbf{Deep Dive: GANITE (\citep{Yoon2018})}

\addcontentsline{toc}{subsubsection}{\protect\numberline{\thesubsubsection}{GANs as Generative Models of Treatment Effect Distributions}}
Although a generative model of the treatment effect distribution is generally unknown, a natural application of GANs is to try to machine learn this distribution from data. GANITE uses two GANs: $GAN_1$, consisting of generator $G_1$ and discriminator $D_1$, to model the counterfactual distribution and $GAN_2$, consisting of generator $G_2$ and discriminator $D_2$, to model the $CATE$ distribution \citep{Yoon2018} (Fig. \ref{fig:ganite}).
The training procedure for $GAN_1$ is as follows:
\begin{enumerate}
\singlespacing
    \item Taking $X$,$T$, and generative noise $Z$ as input, generator $G_1$ generates both potential outcomes $\{\tilde{Y}(T),\tilde{Y}(1-T)\}$. A factual loss $MSE(Y(T),\tilde{Y}(T))$ is applied.
    \item Create a new vector $\textbf{C}=\{Y(T),\tilde{Y}(1-T)\}$ by combining the observed potential outcome and the counterfactual predicted by $G_1$.
    \item Taking $X$ and $\textbf{C}$ as inputs, the discriminator $D_1$ rates each value in $\textbf{C}$ for the probability that it is the observed outcome using the categorical cross entropy loss:
    \begin{equation}
    \mathcal{L}(D_1)=CCE(\{\underbrace{P(\textbf{C}_0=Y(T))}_{\text{Prob first idx is real}},\underbrace{P(\textbf{C}_1=Y(T))}_{\text{Prob sec idx is real}}\},\{\underbrace{\textbf{C}_0==Y(T)}_{\text{1 if idx 0 is real}},\underbrace{\textbf{C}_1==Y(T)}_{\text{1 if idx 1 is real}}\})
\end{equation}
    \item This loss is then fed back to $G_1$ such that the total loss for the generator is now  
    \begin{equation}
        \argmin_{G_1}=MSE(Y(T),\tilde{Y}(T)) - \lambda \mathcal{L}(D_1)
    \end{equation}
\end{enumerate}

After generator $G_1$ is trained to completion, the authors use $\textbf{C}$ as a ``complete dataset" containing both a factual outcome and a counterfactual outcome to train $GAN_2$, which generates treatment effects: 
\begin{enumerate}
\singlespacing
    \item Taking only $X$ and generative noise $Z$ as input, $G_2$ generates a new potential outcome vector $\textbf{R}=\{\hat{Y}(T),\hat{Y}(1-T)\}$. $G_2$ receives an MSE loss to minimize the difference between it's predictions and the ``complete dataset" $\textbf{C}$: $MSE(\textbf{C},\textbf{R})$.
    \item Discriminator $D_2$ takes $X$, $\textbf{C}$, and $\textbf{R}$ as inputs and estimates a probability that $\textbf{C}$ is the ``complete" dataset, and that $\textbf{R}$ is the ``complete dataset":
    \begin{equation}
    \mathcal{L}(D_2)=CCE(\{\underbrace{P(\textbf{C}=\textbf{C})}_{\text{ Prob $\textbf{C}$ is ``CD" }},\underbrace{P(\textbf{R}=\textbf{C})}_{\text{Prob $\textbf{R}$ is ``CD"}}\},\{\underbrace{\textbf{C}==\textbf{C}}_{\text{1 if idx 0 is $\textbf{C}$}},\underbrace{\textbf{C}_1==Y(T)}_{\text{1 if idx 1 is $\textbf{C}$}}\})
\end{equation}
    \item This loss is then fed back to the generator $G_2$ such that the total loss for the generator is now  
    \begin{equation}
        \argmin_{G_2}=MSE(\textbf{C},\textbf{R}) - \lambda \mathcal{L}(D_2)
    \end{equation}
\end{enumerate}
At the end of training, $G_2$ should be able to predict treatment effects with only covariates $X$ and noise $Z$ as inputs. An evolution of GANITE, SCIGAN, extends this framework to settings with more than one treatment and continuous dosages \citep{bica_sicgan}. 

\begin{figure}
    \centering
    \includegraphics[width=\linewidth]{content/GANITE.png}
    \caption{\textbf{GANITE} has two GANs. The first generator $G_1$ generates counterfactuals $\tilde{Y}(T)$. The discriminator $D_1$ attempts to discriminate between these predictions and real data ($Y(T)$). The second generator proposes pairs of potential outcomes $\hat{Y}(0)$ and $\hat{Y}(1)$ (i.e., treatment effects), a vector we call $R$. Discriminator $D_2$ attempts to discern between $R$ and a ``complete dataset" $C$ created by pairing each observed/factual outcome $Y(T)$ with a synthetic outcome $\tilde{Y}(1-T)$ proposed by $G_1$. Although we do not show gradients in other figures, we make an exception for GANs (red line).}
    \label{fig:ganite}
\end{figure}
\newpage
\iffalse
SCIGAN extends GANITE to settings with more than one treatment and continuous dosages \citep{bica_sicgan}. In this setting treatment is parameterized as a binary factor vector denoting which treatment $w \in W$ of possible treatments is used, and a substituent vector of randomly sampled dosages $\mathbb{D}_w$ for treatment $w$, such that $t=(w,\mathbb{D}_w)$. The generator $G$ is identical to $G_1$ in GANITE except that $G$ produces a separate potential outcome for each treatment-dosage combination $\hat{Y}(W,\mathbb{D}_W)=\{\hat{y}(w_1,\mathbb{D}_{w1}),\hat{y}(w_2,\mathbb{D}_{w2}),\dots,\hat{y}(w_n,\mathbb{D}_{wn})\}$. $G$ has two discriminators, $D_w$ and $D_\mathbb{D}$. $D_w$ selects which of the generated treatments $\hat{W}$ is the real one. $D_\mathbb{D}$ subsequently guesses which is the correct dosage for the treatment $\hat{w}$ selected by $D_w$. For out-of-sample prediction, SCIGAN replaces $G_2$ from GANITE with an MLP for out-of-sample prediction using the complete dataset produced by $G$.
\fi
\subsubsection{Adversarial Representation Balancing}
\addcontentsline{toc}{subsubsection}{\protect\numberline{\thesubsubsection}{Adversarial Representation Balancing}}

The use of the IPM loss in CFRNet \citep{Shalit2017} may also be viewed as an adversarial approach in that the representation layers are forced to maximize performance on two competing tasks: predicting outcomes and minimizing an IPM. Rather than using an IPM loss, other authors have trained propensity score estimators that send positive (rather than negative) gradients back to the representation layers \citep{Atan2018,Du2019}. This strategy explicitly decorrelates the covariate representations from the treatment (see Box \ref{box:otherlosses}).

\citet{bica2020estimating} extend this approach to settings with treatment over time using a recurrent neural network. In their medical setting, decorrelating treatment from patient covariates and history allows them to estimate treatment effects at each individual snapshot. This algorithm is described in Section \ref{section:futuredirections}.

\subsubsection{Adversarial IPW Learning}

In a simple adversarial IPW model, \citet{Ozery-flato} estimate IPW weights for the ATE adversarially. A discriminator is presented with two sets of weights: one uniform and the other estimated, and tasked with distinguishing between the two distributions. The ``generator" updates the estimated weights to minimize the ability of the discriminator to distinguish between them. The generator uses exponential gradient ascent during training to regularize the weights and minimize their variance. \citet{KallusDeepMatch} similarly proposes a GAN set-up where a discriminator network attempts to minimize a ``discriminative distance", while the ``generator" of the weights is itself a deep neural network.
